% Options for packages loaded elsewhere
% Options for packages loaded elsewhere
\PassOptionsToPackage{unicode}{hyperref}
\PassOptionsToPackage{hyphens}{url}
\PassOptionsToPackage{dvipsnames,svgnames,x11names}{xcolor}
%
\documentclass[
  brazilian,
  12pt,
]{article}
\usepackage{xcolor}
\usepackage[top=20mm,left=20mm,bottom=20mm,right=20mm]{geometry}
\usepackage{amsmath,amssymb}
\setcounter{secnumdepth}{5}
\usepackage{iftex}
\ifPDFTeX
  \usepackage[T1]{fontenc}
  \usepackage[utf8]{inputenc}
  \usepackage{textcomp} % provide euro and other symbols
\else % if luatex or xetex
  \usepackage{unicode-math} % this also loads fontspec
  \defaultfontfeatures{Scale=MatchLowercase}
  \defaultfontfeatures[\rmfamily]{Ligatures=TeX,Scale=1}
\fi
\usepackage{lmodern}
\ifPDFTeX\else
  % xetex/luatex font selection
  \setmonofont[]{Consolas}
\fi
% Use upquote if available, for straight quotes in verbatim environments
\IfFileExists{upquote.sty}{\usepackage{upquote}}{}
\IfFileExists{microtype.sty}{% use microtype if available
  \usepackage[]{microtype}
  \UseMicrotypeSet[protrusion]{basicmath} % disable protrusion for tt fonts
}{}
\makeatletter
\@ifundefined{KOMAClassName}{% if non-KOMA class
  \IfFileExists{parskip.sty}{%
    \usepackage{parskip}
  }{% else
    \setlength{\parindent}{0pt}
    \setlength{\parskip}{6pt plus 2pt minus 1pt}}
}{% if KOMA class
  \KOMAoptions{parskip=half}}
\makeatother
% Make \paragraph and \subparagraph free-standing
\makeatletter
\ifx\paragraph\undefined\else
  \let\oldparagraph\paragraph
  \renewcommand{\paragraph}{
    \@ifstar
      \xxxParagraphStar
      \xxxParagraphNoStar
  }
  \newcommand{\xxxParagraphStar}[1]{\oldparagraph*{#1}\mbox{}}
  \newcommand{\xxxParagraphNoStar}[1]{\oldparagraph{#1}\mbox{}}
\fi
\ifx\subparagraph\undefined\else
  \let\oldsubparagraph\subparagraph
  \renewcommand{\subparagraph}{
    \@ifstar
      \xxxSubParagraphStar
      \xxxSubParagraphNoStar
  }
  \newcommand{\xxxSubParagraphStar}[1]{\oldsubparagraph*{#1}\mbox{}}
  \newcommand{\xxxSubParagraphNoStar}[1]{\oldsubparagraph{#1}\mbox{}}
\fi
\makeatother


\usepackage{longtable,booktabs,array}
\usepackage{calc} % for calculating minipage widths
% Correct order of tables after \paragraph or \subparagraph
\usepackage{etoolbox}
\makeatletter
\patchcmd\longtable{\par}{\if@noskipsec\mbox{}\fi\par}{}{}
\makeatother
% Allow footnotes in longtable head/foot
\IfFileExists{footnotehyper.sty}{\usepackage{footnotehyper}}{\usepackage{footnote}}
\makesavenoteenv{longtable}
\usepackage{graphicx}
\makeatletter
\newsavebox\pandoc@box
\newcommand*\pandocbounded[1]{% scales image to fit in text height/width
  \sbox\pandoc@box{#1}%
  \Gscale@div\@tempa{\textheight}{\dimexpr\ht\pandoc@box+\dp\pandoc@box\relax}%
  \Gscale@div\@tempb{\linewidth}{\wd\pandoc@box}%
  \ifdim\@tempb\p@<\@tempa\p@\let\@tempa\@tempb\fi% select the smaller of both
  \ifdim\@tempa\p@<\p@\scalebox{\@tempa}{\usebox\pandoc@box}%
  \else\usebox{\pandoc@box}%
  \fi%
}
% Set default figure placement to htbp
\def\fps@figure{htbp}
\makeatother


% definitions for citeproc citations
\NewDocumentCommand\citeproctext{}{}
\NewDocumentCommand\citeproc{mm}{%
  \begingroup\def\citeproctext{#2}\cite{#1}\endgroup}
\makeatletter
 % allow citations to break across lines
 \let\@cite@ofmt\@firstofone
 % avoid brackets around text for \cite:
 \def\@biblabel#1{}
 \def\@cite#1#2{{#1\if@tempswa , #2\fi}}
\makeatother
\newlength{\cslhangindent}
\setlength{\cslhangindent}{1.5em}
\newlength{\csllabelwidth}
\setlength{\csllabelwidth}{3em}
\newenvironment{CSLReferences}[2] % #1 hanging-indent, #2 entry-spacing
 {\begin{list}{}{%
  \setlength{\itemindent}{0pt}
  \setlength{\leftmargin}{0pt}
  \setlength{\parsep}{0pt}
  % turn on hanging indent if param 1 is 1
  \ifodd #1
   \setlength{\leftmargin}{\cslhangindent}
   \setlength{\itemindent}{-1\cslhangindent}
  \fi
  % set entry spacing
  \setlength{\itemsep}{#2\baselineskip}}}
 {\end{list}}
\usepackage{calc}
\newcommand{\CSLBlock}[1]{\hfill\break\parbox[t]{\linewidth}{\strut\ignorespaces#1\strut}}
\newcommand{\CSLLeftMargin}[1]{\parbox[t]{\csllabelwidth}{\strut#1\strut}}
\newcommand{\CSLRightInline}[1]{\parbox[t]{\linewidth - \csllabelwidth}{\strut#1\strut}}
\newcommand{\CSLIndent}[1]{\hspace{\cslhangindent}#1}

\ifLuaTeX
\usepackage[bidi=basic]{babel}
\else
\usepackage[bidi=default]{babel}
\fi
% get rid of language-specific shorthands (see #6817):
\let\LanguageShortHands\languageshorthands
\def\languageshorthands#1{}


\setlength{\emergencystretch}{3em} % prevent overfull lines

\providecommand{\tightlist}{%
  \setlength{\itemsep}{0pt}\setlength{\parskip}{0pt}}



 


\usepackage{authblk}
\makeatletter
\@ifpackageloaded{float}{}{\usepackage{float}}
\floatstyle{plain}
\@ifundefined{c@chapter}{\newfloat{algo}{h}{loalgo}}{\newfloat{algo}{h}{loalgo}[chapter]}
\floatname{algo}{Algoritmo}
\newcommand*\listofalgos{\listof{algo}{List of Algoritmos}}
\makeatother
\makeatletter
\@ifpackageloaded{caption}{}{\usepackage{caption}}
\AtBeginDocument{%
\ifdefined\contentsname
  \renewcommand*\contentsname{Índice}
\else
  \newcommand\contentsname{Índice}
\fi
\ifdefined\listfigurename
  \renewcommand*\listfigurename{Lista de Figuras}
\else
  \newcommand\listfigurename{Lista de Figuras}
\fi
\ifdefined\listtablename
  \renewcommand*\listtablename{Lista de Tabelas}
\else
  \newcommand\listtablename{Lista de Tabelas}
\fi
\ifdefined\figurename
  \renewcommand*\figurename{Figura}
\else
  \newcommand\figurename{Figura}
\fi
\ifdefined\tablename
  \renewcommand*\tablename{Tabela}
\else
  \newcommand\tablename{Tabela}
\fi
}
\@ifpackageloaded{float}{}{\usepackage{float}}
\floatstyle{ruled}
\@ifundefined{c@chapter}{\newfloat{codelisting}{h}{lop}}{\newfloat{codelisting}{h}{lop}[chapter]}
\floatname{codelisting}{Listagem}
\newcommand*\listoflistings{\listof{codelisting}{Lista de Listagens}}
\usepackage{amsthm}
\theoremstyle{definition}
\newtheorem{example}{Exemplo}[section]
\theoremstyle{definition}
\newtheorem{definition}{Definição}[section]
\theoremstyle{remark}
\AtBeginDocument{\renewcommand*{\proofname}{Comprovação}}
\newtheorem*{remark}{Comentário}
\newtheorem*{solution}{Solução}
\newtheorem{refremark}{Comentário}[section]
\newtheorem{refsolution}{Solução}[section]
\makeatother
\makeatletter
\makeatother
\makeatletter
\@ifpackageloaded{caption}{}{\usepackage{caption}}
\@ifpackageloaded{subcaption}{}{\usepackage{subcaption}}
\makeatother
\makeatletter
\@ifpackageloaded{algorithm}{}{\usepackage{algorithm}}
\makeatother
\makeatletter
\@ifpackageloaded{algpseudocode}{}{\usepackage{algpseudocode}}
\makeatother
\makeatletter
\@ifpackageloaded{caption}{}{\usepackage{caption}}
\makeatother
    \makeatletter
    \newcommand\fs@nocaption{
      \def\@fs@cfont{\bfseries}
      \let\@fs@capt\floatc@plain
      \def\@fs@pre{}%
      \def\@fs@post{\kern2pt\hrule}%
      \def\@fs@mid{\hrule\kern2pt}%
      \let\@fs@iftopcapt\iftrue}
    \makeatother
  
\captionsetup[algorithm]{justification=raggedright}
\usepackage{bookmark}
\IfFileExists{xurl.sty}{\usepackage{xurl}}{} % add URL line breaks if available
\urlstyle{same}
\hypersetup{
  pdftitle={Aproximações para a duração média de passeios aleatórios},
  pdflang={pt-BR},
  pdfkeywords={Passeio Aleatório, Monte Carlo, Estimação},
  colorlinks=true,
  linkcolor={blue},
  filecolor={Maroon},
  citecolor={Blue},
  urlcolor={Blue},
  pdfcreator={LaTeX via pandoc}}


\title{Aproximações para a duração média de passeios aleatórios}
\author{Eduardo Yukio Garrafa Ishihara \and Gustavo Silva
Garone \and Elisabeti Kira}
\date{10 de fevereiro de 2026}
\begin{document}
\maketitle
\begin{abstract}
Inspirados na física, propusemos aproximações para a duração de passeios
aleatórios diversos, em especial, os de difícil derivação analítica.
Para avaliar as aproximações, utilizamos simulações de Monte Carlo na
linguagem Julia. Analisamos o desempenho desta linguagem comparando-a
com a linguagem Python. Com modelos (passeios) mais sofisticados,
utilizamos recursos computacionais para solução de sistemas lineares na
construção de novas aproximações, descrevendo o uso de paralelismo e de
outras ferramentas. Exibimos e analisamos os resultados obtidos, assim
como os algoritmos utilizados.
\end{abstract}

\renewcommand{\Return}{\State \textbf{return}~}
\newcommand{\Print}{\State \textbf{print}~}
\newcommand{\Break}{\State \textbf{break}}
\newcommand{\Continue}{\State \textbf{continue}}
\newcommand{\True}{\textbf{true}}
\newcommand{\False}{\textbf{false}}
\renewcommand{\And}{\textbf{and}~}
\newcommand{\Or}{\textbf{or}~}
\renewcommand{\Not}{\textbf{not}~}
\newcommand{\To}{\textbf{to}~}
\newcommand{\DownTo}{\textbf{downto}~}

\floatname{algorithm}{Algoritmo}


\textsubscript{Fonte:
\href{https://GustavoGarone.github.io/Estimador-de-Passeio/Estimando_Passeios_Aleatorios.qmd.html}{Caderno
de Artigo}}

\section{Introdução}\label{introduuxe7uxe3o}

Neste artigo, estudaremos a duração de passeios aleatórios e suas
aproximações. Temos como objetivo descrever modelos de passeios como a
``Ruína do Jogador'' de Pascal, discutido em EDWARDS
(\citeproc{ref-edwards_pascals_1983}{1983}), generalizando-o para
múltiplas dimensões. Abordaremos técnicas analíticas, computacionais e
de estimativas para a duração esperada destes passeios.

Começaremos por definir o modelo que será utilizado para descrever os
passeios aleatórios, acompanhado de exemplo. Em seguida, descreveremos
resultados anteriores do cálculo de probabilidades envolvendo esses
passeios. Então, apresentaremos um estimador para a duração desses e
partiremos para sua avaliação. Compará-lo-emos com os resultados
probabilísticos obtidos. Quando avançarmos para passeios em que
resultados analíticos não são factíveis, iniciaremos a comparação do
estimador com simulações de Monte Carlo.

Descrevemos e justificamos pela literatura a adoção da linguagem de
programação \emph{Julia}, essencial para possibilitar a execução
eficiente de simulações computacionalmente intensivas em passeios mais
complexos.

Em dimensões maiores, forneceremos, até onde é computacionalmente
possível, resultados exatos para a duração esperada de passeios
aleatórios. Para obter esses resultados, utilizaremos métodos para
solução de sistemas lineares extensos e empregaremos técnicas de
computação paralela, que serão brevemente discutidas.

Este artigo cobrirá tanto o desenvolvimento teórico matemático do
estimador e modelos, quanto a descrição do aspecto computacional desta
pesquisa, incluindo os algoritmos utilizados. Os códigos completos estão
disponíveis para consulta e \emph{Download} na versão digital desta
publicação.

\section{A Ruína do Jogador e outros passeios
aleatórios}\label{a-ruuxedna-do-jogador-e-outros-passeios-aleatuxf3rios}

Postulado por Blaise Pascal para Fermat no século XVII
(\citeproc{ref-edwards_pascal_1983}{\textbf{edwards\_pascal\_1983?}}), o
problema que originaria a Ruína do Jogador consistia em determinar a
probabilidade de vitória de um jogador em um jogo de apostas, como o
lançamento de dados ou moedas, contra um adversário. A característica
principal deste jogo é que os jogadores transferem seus pontos, ou
fortunas, entre si. Portanto, o jogo termina quando um dos jogadores
perde toda a sua fortuna, ou seja, vai à ruína.

Podemos descrever a fortuna do jogador nesses jogos como um passeio
aleatório unidimensional em \(\mathbb{Z}\), usualmente considerando
apenas o espaço de estados \(\mathbb{N}\), com as probabilidades
envolvidas nas jogadas controlando o comportamento do passeio. Por
exemplo, na aposta sobre o resultado do lançamento de uma moeda honesta,
temos um passeio simétrico:

\[
\begin{aligned}
\mathbb{P}(X = +1) &= 0.5, \\
\mathbb{P}(X = -1) &= 0.5.
\end{aligned}
\]

A ``Ruína do Jogador'' é o nome do resultado que, num jogo contra uma
``casa de apostas'' com fortuna infinita ou muito maior que a sua, o
jogador inevitavelmente irá à ruína quando o jogo for justo (quando o
passeio configurar um martingal) em tempo finito.

Podemos estender esse passeio para jogos mais complexos, desde
alterações simples, como probabilidades desiguais, até jogos com mais
jogadores, fortunas distintas (múltiplas dimensões) e regras variáveis
ao longo do tempo (jogada) ou espaço (fortuna). Assim, conseguimos
modelar uma grande variedade de passeios aleatórios com base no
princípio da troca de fortuna.

\section{Modelando passeios
aleatórios}\label{modelando-passeios-aleatuxf3rios}

\subsection{Modelo}\label{modelo}

\begin{definition}[]\protect\hypertarget{def-modelodiscreto}{}\label{def-modelodiscreto}

Considere que o jogador \(A\) começa com uma fortuna inicial com um
total \(\$N\) em jogo, \(\$a\), \(0 \leq a \leq N\), fornecendo um
espaço de estados \(\{0, 1, \dots, a, \dots, N\}\). Cada rodada produz
um resultado aleatório independente \(X_i\), \(i=1,2,\dots\),
representando o ganho ou perda naquela etapa. Isso será chamado de
variável ``regra''. Após \(n\) rodadas, a fortuna do jogador \(A\) é
dada por

\[
S_n = a + \sum_{i=1}^n X_i.
\]

A \textbf{duração do jogo} é o primeiro tempo \(n\) tal que
\(S_n \leq 0\) (ruína) ou \(S_n \geq N\) (vitória). É denotada por
\(\tau_{X, a, N}\), a regra de parada.

Podemos também relacionar \(X\) com a posição atual do passeio,
\(X_{n} := X_{S_n}\) estendendo esse modelo para passeios em meio não
homogêneo: definimos \(\pmb{R} =
\{X_1, \dots, X_N\}\) e \(\tau_{\pmb{R}, a, N}\).

\end{definition}

\begin{refremark}
Neste artigo, não trataremos de passeios em meio não homogêneo no tempo,
mas podemos, se preciso, estender esse modelo para incluir essa
característica.

\label{rem-tempo}

\end{refremark}

Para ajudar no entendimento deste modelo, oferecemos um exemplo:

\begin{example}[]\protect\hypertarget{exm-modelodiscreto}{}\label{exm-modelodiscreto}

Considere um jogador apostando em um casino no lançamento de uma moeda
honesta. Se a moeda cair em cara, o jogador ganha um dólar; se cair em
coroa, perde um dólar. O jogo termina quando o jogador atingir cem
dólares ou vá à ruína (seu saldo chegue em zero). O jogador começa com
cinquenta dólares.

Podemos modelar esse jogo no modelo da
Definição~\ref{def-modelodiscreto}:

Seja \(a = 50\) e \(N = 100\). A variável regra do jogo é dada por

\[
X = \begin{cases}
-1, & \mathbb{P}(X = -1) = 0.4, \\
+1, & \mathbb{P}(X = +1) = 0.6.
\end{cases}
\]

Nota-se que a regra é a mesma para todo estado \(S_n\) enquanto o jogo
durar, ou seja, o passeio está homogêneo no espaço e tempo.

Após \(5\) jogadas, o jogo tomou o seguinte percurso: \[
\pmb{x} = \{x_1, x_2, x_3, x_4, x_5\} = \{1, -1, 1, 1, -1\}.
\]

temos então que a posição (fortuna) atual do jogador é \[
S_5 = 50 + (1 - 1 + 1 + 1 -1) = 51.
\]

\end{example}

\subsection{Jogos unitários e não
unitários}\label{jogos-unituxe1rios-e-nuxe3o-unituxe1rios}

Definimos um jogo como \textbf{unitário} se a variável regra \(X\)
assume apenas valores inteiros unitários, ou seja,
\(X: \Omega \rightarrow \{-1, +1\}\). Caso contrário, o jogo é dito
\textbf{não unitário}. No Exemplo~\ref{exm-modelodiscreto}, o jogo é
unitário.

Essa distinção faz-se importante, uma vez que, como apresentaremos na
próxima seção, a abordagem analítica para determinar a duração esperada
do jogo difere entre jogos unitários e não unitários. De fato, em jogos
não unitários, a determinação analítica de
\(\mathbb{E}[\tau_{X, a, N}]\) pode ser muito mais complexa, ou até
mesmo inviável.

\subsection{\texorpdfstring{Abordagem analítica para determinar
\(\mathbb{E}[\tau_{X, a, N}]\)}{Abordagem analítica para determinar \textbackslash mathbb\{E\}{[}\textbackslash tau\_\{X, a, N\}{]}}}\label{sec-anali}

Retomamos a variável aleatória de regra \(X_i, i \in \mathbb{Z}^+\),
para descrever a quantia ganha pelo jogador \(A\) na \(i\)-ésima rodada,
Podemos considerar \(X_i, i \in \mathbb{Z}^+\), independentes e com
mesma distribuição de \(X\) dada por

\[
X = \begin{cases}
    +1, & \mathbb{P}(X=+1)=p \\
    -1, & \mathbb{P}(X=-1)=q
\end{cases}
\]

Adicionalmente, definimos \(\tau_{X,a,N}\), como o número de rodadas até
o fim do jogo em que o resultado de cada rodada é representado pela
variável aleatória \(X\), o jogador \(A\) começa com \(a\) reais e o
jogo termina quando \(A\) chega a \(0\) ou \(N\) reais. Definimos
\(\mu_{a} = \mathbb{E}[\tau_{X,a,N}]\), ou seja, é o número esperado de
rodadas que durará o jogo em que que o resultado de cada rodada é
representado pela variável aleatória \(X\), o jogador \(A\) começa com
\(a\) reais e o jogo termina quando \(A\) chega a \(0\) ou \(N\) reais.

\[
\begin{aligned}
\mathbb{E}[\tau_{X, a, N}]
&= \mathbb{E}[\tau_{X, a, N}|X_1=1] \mathbb{P}(X_1=1)+\mathbb{E}[\tau_{X, a, N}|X_1=-1] \mathbb{P}(X_1=-1) \\
&= \mathbb{E}[\tau_{X, a, N}|X_1=1] \: p + \mathbb{E}[\tau_{X, a, N}|X_1=-1] \: q \\
&= (\mathbb{E}[\tau_{X, a+1,N}]+1) \: p + (\mathbb{E}[\tau_{X, a-1,N}]+1) \: q \\
&= p \: \mathbb{E}[\tau_{X, a+1,N}] + q \: \mathbb{E}[\tau_{X, a-1,N}] + 1 \\
\mu_{a} &= p \mu_{a+1} + q \mu_{a-1} + 1  
\end{aligned}
\] Trata-se, portanto, de uma equação de diferenças finitas não
homogênea, cuja solução é a esperança desejada.

Tomemos \(\mu_{a} = \lambda ^a\) e, temporariamente, desconsideraremos a
constante \(1\) na equação de diferenças finitas. \[
\begin{aligned}
& \phantom{\Rightarrow;;}\lambda ^a = p \:  \lambda ^{a+1} + q \: \lambda ^{a-1} \\
& \Rightarrow p \:  \lambda ^ 2 -\lambda + q = 0 \\
& \Rightarrow \lambda_1 = \frac{q}{p}; \:  \lambda_2 = 1
\end{aligned}
\]

Para \(\frac{p}{q} = 1\): \[
\lambda_1 = \lambda_2 = 1
\] Temos a seguinte equação característica: \[
\mu_{a} = \alpha + \beta a + \gamma a^2
\]

Tomemos \(\mu_{a} = \gamma a^2\). Juntando as equações: \[
\begin{aligned}
& \phantom{\Rightarrow;;}
\gamma a^2 = \frac{1}{2} \gamma (a-1)^2 + \frac{1}{2} \gamma (a+1)^2 + 1 \\
& \Rightarrow 2 \gamma a^2 = \gamma a^2 - 2 \gamma a + \gamma + \gamma a^2 + 2\gamma a + \gamma + 2\\
& \Rightarrow \gamma = -1
\end{aligned}
\]

Aplicando as seguintes condições de fronteira: \[
\begin{aligned}
& \phantom{\Rightarrow \,\,}\begin{cases} \mu_{0} = 0 \\ \mu_{N} = 0 \end{cases} \\
&\Rightarrow
\begin{cases} \mu_{0} = 0 = \alpha + \beta 0 + \gamma 0^2 \\ 
              \mu_{N} = 0 = \alpha + \beta N + \gamma N^2   \end{cases}\\
& \Rightarrow 
\begin{cases} \alpha = 0 \\ 
              \beta  = N  \end{cases} \\
\mu_{a} &=  Na-a^2 = a(N-a)
\end{aligned}
\] Nota-se que, neste modelo, a esperança da duração do jogo é máxima
quando os dois jogadores começam com a mesma fortuna.

Para jogos não simétricos, isto é, jogos em que
\(\mathbb{E}[X] \neq 0\), temos a seguinte equação característica: \[
\mu_{a} = \alpha \lambda_1^a + \beta \lambda_2^a + \gamma a
\]

Tomemos \(\mu_{a} = \gamma a\):

\[
\begin{aligned}
& \phantom{\Rightarrow;;}\gamma a = p \:  \gamma ({a+1}) + q \: \gamma ({a-1}) +1 \\
& \Rightarrow \gamma = \frac{1}{q-p}
\end{aligned}
\]

Substituindo e usando as fronteiras: \[
\begin{aligned}
& \phantom{\Rightarrow;;} \mu_{a} = \alpha \left( \frac{q}{p} \right) ^a + \beta + \frac{a}{q-p} \\
& \Rightarrow \begin{cases} \mu_{0} = 0 \\ \mu_{N} = 0 \end{cases}\\
& \Rightarrow
\begin{cases} \mu_{0} = 0 = \alpha \left( \frac{q}{p} \right) ^0 + \beta + \frac{0}{q-p} \\ 
              \mu_{N} = 0 = \alpha \left( \frac{q}{p} \right) ^N + \beta + \frac{N}{q-p}   \end{cases}\\
& \Rightarrow 
\begin{cases} \alpha = \frac{N}{p-q}\frac{1}{\left( \frac{q}{p} \right)^N-1} \\ 
              \beta  = -\frac{N}{p-q}\frac{1}{\left( \frac{q}{p} \right)^N-1} \end{cases} \\ \\ \\
\mu_{a} &= \frac{N}{p-q}\frac{1}{\left( \frac{q}{p} \right)^N-1} \left( \frac{q}{p} \right)^a - \frac{N}{p-q}\frac{1}{\left( \frac{q}{p} \right)^N-1} + \frac{a}{q-p} \\
\mu_{a} & = \frac{a}{q-p}-\frac{N}{q-p} \left( \frac{1- \left( \frac{q}{p} \right)^a}
        {1- \left( \frac{q}{p} \right)^N}\right)
\end{aligned}
\]

Finalmente, \[
\mathbb{E}[\tau_{x, a, N}] = 
\begin{cases}
a(N-a), & \mathbb{E}[X] = 0 \\
\frac{a}{q-p}-\frac{N}{q-p} \left( \frac{1- \left( \frac{q}{p} \right)^a} {1- \left( \frac{q}{p} \right)^N}\right), &
\mathbb{E}[X] \neq 0.
\end{cases}
\]

O resultado final é o mesmo que o obtido por outros métodos STERN
(\citeproc{ref-stern_conditional_1975}{1975}).

\section{Aproximando durações de passeios
aleatórios}\label{aproximando-durauxe7uxf5es-de-passeios-aleatuxf3rios}

Uma pergunta natural quando se discute aproximações é quanto a sua
utilidade. Se possível, é mais desejável utilizar resultados obtidos
analiticamente, como, no exemplo da ``Ruína do Jogador'' clássica, os
apresentados na Seção~\ref{sec-anali}.

Nos cenários em que não é prático o desenvolvimento analítico, soluções
computacionais são buscadas. Estas soluções podem fornecer resultados
teóricos corretos, como na solução de sistemas lineares extensos que
abordaremos na {[}SESSÃO AQUI{]}, ou aproximações pelo método de Monte
Carlo, como descreve HARRISON
(\citeproc{ref-harrison_introduction_2010}{2010}).

Ainda assim, a simulação também traz desvantagens para além da perda de
precisão e de interpretabilidade. Por exemplo, como abordado em RITTER
et al. (\citeproc{ref-ritter_determining_2011}{2011}), um número
considerável de simulações é necessário para garantir acurácia dos dados
e estabilidade. Diante disso, vê-se justificável a busca de estimadores
eficientes para a duração esperada de passeios aleatórios diversos.

\subsection{Um estimador heurístico}\label{um-estimador-heuruxedstico}

\begin{definition}[]\protect\hypertarget{def-estimadorsimples}{}\label{def-estimadorsimples}

Seja um jogo definido conforme a Definição~\ref{def-modelodiscreto}, com
variável regra \(X\) tal que \(\mathbb{E}(X) \neq 0\), valores \(a\) e
\(N\). O estimador proposto para a duração esperada desse jogo é dado
por:

\[
\hat{\mathbb{E}}(\tau_{X, a, N}) = \begin{cases}
\frac{N-a}{\mathbb{E}(X)},  & \mathbb{E}(X) > 0, \\
\left \lvert \frac{a}{\mathbb{E}(X)} \right \rvert,  & \mathbb{E}(X) < 0.
\end{cases}
\]

\end{definition}

Primeiramente, o estimador da Definição~\ref{def-estimadorsimples} foi
comparado com o resultado da \textbf{?@eq-tempoteorico} para jogos
unitários. Retomando o Exemplo~\ref{exm-modelodiscreto}, temos, como
\(\mathbb{E}(X) = 0.2\):

\[
\begin{aligned}
\hat{\mathbb{E}}(\tau_{X, 50, 100}) = \frac{100}{0.2} = 250 \\
\mathbb{E}(\tau_{X, 50, 100}) = \frac{50}{0.6 - 0.4} - \frac{100}{0.6 - 0.4}
\left( \frac{1 - \left( \frac{0.4}{0.6} \right)^{50}}{1 - \left( \frac{0.4}{0.6}
\right)^{100}} \right) = 250
\end{aligned}
\]

Partimos para testá-lo em outros casos unitários. Dos nossos resultados,
obtemos, fixada a fortuna inicial \(a=50\) e total \(N\), e variando a
probabilidade de ganho \(p\):

\textsubscript{Fonte:
\href{https://GustavoGarone.github.io/Estimador-de-Passeio/Estimando_Passeios_Aleatorios.qmd.html}{Caderno
de Artigo}}

\begin{figure}[H]

\centering{

\pandocbounded{\includegraphics[keepaspectratio]{Estimando_Passeios_Aleatorios_files/mediabag/Estimando_Passeios_Aleatorios_files/figure-pdf/fig-comparacaoestimadorxteorico-output-1.pdf}}

}

\caption{\label{fig-comparacaoestimadorxteorico}Comparação entre
estimador simples e duração esperado em jogos unitários para \(a = 50\)
e \(N = 100\) em função de \(p\).}

\end{figure}%

\textsubscript{Fonte:
\href{https://GustavoGarone.github.io/Estimador-de-Passeio/Estimando_Passeios_Aleatorios.qmd.html}{Caderno
de Artigo}}

Notamos que o estimador apresenta bom comportamento nessa classe de
jogos. Uma inadequação de seu uso existe para casos em que
\(\mathbb{E}(X)\) se aproxima de \(0\), como indica o gráfico da
Figura~\ref{fig-comparacaoestimadorxteorico} no intervalo
\([0.48, 0.52]\)

Validado o estimador para grande parte dos casos mais simples, partimos
para análise em jogos não unitários. Para isso, utilizamos simulações de
Monte Carlo para obter valores aproximados de
\(\mathbb{E}(\tau_{X, a, N})\). Apresentamos no
Algoritmo~\ref{algo-montecarlosimples} uma implementação simples desse
processo HARRISON (\citeproc{ref-harrison_introduction_2010}{2010}),
RITTER et al. (\citeproc{ref-ritter_determining_2011}{2011}).

\floatstyle{ruled}
\restylefloat{algorithm}
\begin{algorithm}[htb!]
\caption{\label{algo-montecarlosimples}Monte Carlo Simples}
\begin{algorithmic}[1]
\Procedure{MonteCarlo}{$MC, p, g, l, a, N$}
    \State $total\_tempo \gets 0$
    \For{$i =1$ to $MC$}
        \State $pos$ \gets $a$
        \While{$pos > 0$ and $pos < N$}
            \State $u \gets \Call{Sortear}{U(0,1)}$
            \If{$u < p$}
                \State $pos$ \gets $pos + g$
            \Else
                \State $pos$ \gets $pos - l$
            \EndIf
            \State $total\_tempo \gets total\_tempo + 1$
        \EndWhile
    \EndFor
    \State \Return $total\_tempo / MC$
\EndProcedure
\end{algorithmic}
\end{algorithm}
\floatstyle{plain}

\begin{figure}[H]

\centering{

\pandocbounded{\includegraphics[keepaspectratio]{Estimando_Passeios_Aleatorios_files/mediabag/Estimando_Passeios_Aleatorios_files/figure-pdf/fig-comparacaoestimadorxsimulado-output-1.pdf}}

}

\caption{\label{fig-comparacaoestimadorxsimulado}Comparação entre
estimador simples e duração simulada para \(a = 50\), \(N = 100\), valor
de perda \$1 e de ganho \$2, em função de \(p\). Simulados com
\(10.000\) iterações de Monte Carlo.}

\end{figure}%

\textsubscript{Fonte:
\href{https://GustavoGarone.github.io/Estimador-de-Passeio/Estimando_Passeios_Aleatorios.qmd.html}{Caderno
de Artigo}}

O estimador também apresenta bom desempenho para jogos não unitários,
como mostra a Figura~\ref{fig-comparacaoestimadorxsimulado}. Novamente,
observa-se uma queda na precisão do estimador à medida que
\(\mathbb{E}(X)\) se aproxima de \(0\).

Com esses resultados, validamos o estimador proposto na
Definição~\ref{def-estimadorsimples} para passeios aleatórios em meio
homogêneo, tanto unitários quanto não unitários. Dessa forma, partimos
para a análise de passeios mais complexos.

\subsection{Meio não-homogêneo}\label{meio-nuxe3o-homoguxeaneo}

Ao tratarmos de passeios em meio não-homogêneo, precisamos descrever o
comportamento deste ambiente. No geral, podemos modelar um meio
aleatório ou determinístico. Em ambos casos, a variável regra \(X\)
dependerá da posição atual do passeio, ou seja, \(X = X_{S_n}\).
Entretanto, enquanto no meio aleatório pode ser modelado

\newpage{}

\section{Referências}\label{referuxeancias}

\protect\phantomsection\label{refs}
\begin{CSLReferences}{0}{1}
\bibitem[\citeproctext]{ref-edwards_pascals_1983}
EDWARDS, A. W. F. \href{https://doi.org/10.2307/1402732}{Pascal's
Problem: The 'Gambler's Ruin'}. \textbf{International Statistical Review
/ Revue Internationale de Statistique}, v. 51, n. 1, p. 73--79, 1983.

\bibitem[\citeproctext]{ref-harrison_introduction_2010}
HARRISON, R. L. \href{https://doi.org/10.1063/1.3295638}{Introduction To
Monte Carlo Simulation}. \textbf{{AIP} conference proceedings}, v. 1204,
p. 17--21, 5 jan. 2010.

\bibitem[\citeproctext]{ref-ritter_determining_2011}
RITTER, F. E. et al.
\href{https://doi.org/10.1007/978-0-85729-883-6_5}{Determining the
Number of Simulation Runs: Treating Simulations as Theories by Not
Sampling Their Behavior}. Em: ROTHROCK, L.; NARAYANAN, S. (Eds.).
\textbf{Human-in-the-Loop Simulations: Methods and Practice}. London:
Springer, 2011. p. 97--116.

\bibitem[\citeproctext]{ref-stern_conditional_1975}
STERN, F.
\href{https://doi.org/10.1080/0025570X.1975.11976486}{Conditional
Expectation of the Duration in the Classical Ruin Problem}.
\textbf{Mathematics Magazine}, v. 48, n. 4, p. 200--203, 1 set. 1975.

\end{CSLReferences}




\end{document}
